%! Advanced Factoring — Unit 3, Lesson 3.2 — Warm-Up
\documentclass[10pt]{article}
\usepackage{algebra2-article}
\usepackage{algebra2-boxes}
\newcommand{\TallMath}[1]{$\displaystyle #1\rule[-1.4em]{0pt}{3.2em}$}

\begin{document}

\pageheader{Unit 3, Lesson 3.2}{Warm-Up}
\namedateperiod

\vspace{0.10in}

\begin{notesbox}{Getting Ready: Skills We'll Use Today}
\small Yesterday you multiplied polynomials \emph{forward}. Today you will run them \emph{backward}.
These four skills are what that takes.

\vspace{5pt}
\textbf{1. Greatest common factor.} The GCF of two monomials uses the largest common number and the
\emph{lowest} power of each shared variable.

\vspace{2pt}
GCF of $18$ and $24$: \blank{1.2cm} \qquad GCF of $12x^3$ and $18x^5$: \blank{1.8cm} \qquad
GCF of $10x^2y$ and $15xy^3$: \blank{2.0cm}

\vspace{3pt}
Factor it out: \; $10x^2 + 15x = $ \blank{4.0cm}

\vspace{6pt}
\textbf{2. Factoring you already own} (from Lesson 2.2).

\vspace{2pt}
$x^2 + 7x + 12 = $ \blank{3.6cm} \qquad $x^2 - 9 = $ \blank{3.6cm}

\vspace{3pt}
Which one used the \emph{difference of squares} pattern $a^2-b^2=(a+b)(a-b)$? \blank{2.6cm}

\vspace{6pt}
\textbf{3. A product from yesterday --- watch what survives.} Multiply, then collect like terms.

\vspace{2pt}
$(x + 2)(x^2 - 2x + 4)$

\vspace{2pt}
Six products: \blank{1.5cm} \, \blank{1.5cm} \, \blank{1.5cm} \, \blank{1.5cm} \, \blank{1.5cm} \,
\blank{1.5cm}

\vspace{3pt}
Collected answer: \blank{3.6cm} \quad How many of the six terms survived? \blank{1.2cm}

\vspace{6pt}
\textbf{4. Perfect cubes.} Fill the table --- you will need these on sight today.

\vspace{2pt}
\begin{center}
\renewcommand{\arraystretch}{1.4}
\begin{tabular}{|c|c|c|c|c|c|}
\hline
$n$ & $1$ & $2$ & $3$ & $4$ & $5$ \\ \hline
$n^3$ & \blank{0.9cm} & \blank{0.9cm} & \blank{0.9cm} & \blank{0.9cm} & \blank{0.9cm} \\ \hline
\end{tabular}
\end{center}

\vspace{2pt}
Also: $(2x)^3 = $ \blank{1.6cm} \qquad $(3y)^3 = $ \blank{1.6cm}

\end{notesbox}

\end{document}
